\documentclass[11pt]{article}
% Shared preamble for per-Python-file documentation (input via \input{...}).
\usepackage[margin=1in]{geometry}
\usepackage[T1]{fontenc}
\usepackage[utf8]{inputenc}
\usepackage{lmodern}
\usepackage{hyperref}
\usepackage{xcolor}
\usepackage{listings}
\usepackage{listingsutf8}
\lstset{
  basicstyle=\ttfamily\small,
  columns=fullflexible,
  breaklines=true,
  upquote=true,
  inputencoding=utf8,
  extendedchars=true,
  frame=single,
  framerule=0.2pt,
  tabsize=2
}


\title{\texttt{model/egmn.py} (Q\&A)}
\author{}
\date{}

\begin{document}
\maketitle

\paragraph{Q: What is the purpose of \texttt{model/egmn.py}?}
\textbf{A:} It implements the proposed \textbf{EGMN} (Exponential--Gaussian Mixture Network) model for watch-time prediction. The model produces mixture parameters and provides both a likelihood-based training loss and a mean-based point predictor.

\paragraph{Q: What is the model input and output interface?}
\textbf{A:} The forward pass \lstinline|forward(x_dict)| takes a feature dictionary (as produced by \lstinline|dataloader/kuairec.py|) and returns \lstinline|(pi, lambda_, mu, sigma)|. The \lstinline|predict(x)| method returns a scalar point prediction per example computed as a mixture mean.

\paragraph{Q: How does EGMN encode features?}
\textbf{A:} It builds a feature registry from the dataset \lstinline|description|, uses \lstinline|Embedding| layers for sparse and sequence features, uses linear layers for continuous features, averages masked sequence embeddings, concatenates all representations, and feeds them into a shared MLP trunk.

\paragraph{Q: What mixture components does EGMN parameterize?}
\textbf{A:} The code uses one Exponential component (rate \lstinline|lambda_|) and multiple Normal components (means \lstinline|mu| and standard deviations \lstinline|sigma|). Mixture weights are produced from \lstinline|pi| logits via softmax.

\paragraph{Q: What does \lstinline|loss(...)| compute?}
\textbf{A:} It computes (1) a negative log-likelihood term via log-sum-exp over mixture components, (2) an L1 reconstruction term between the label and the mixture-mean prediction, and (3) an entropy-style regularizer on mixture weights. The run script combines these with weights \lstinline|alpha| and \lstinline|beta|.

\paragraph{Q: Code example: mixture-mean prediction used in evaluation?}
\textbf{A:}
\begin{lstlisting}[language=Python]
pi, lambda_, mu, sigma = model(features)
weights = torch.softmax(pi, dim=1)
component_means = torch.concat([1 / lambda_, mu], dim=1)
y_hat = torch.sum(weights * component_means, dim=1)
\end{lstlisting}

\paragraph{Q: Full source code?}
\textbf{A:}
\lstinputlisting[language=Python]{/workspace/model/egmn.py}

\end{document}
